\documentclass{tufte-handout}

\usepackage{style}

\begin{document}

\tma{04} % TMA number

\begin{question}

Let \( a \) be the \( y \)-intercept and \( \theta \) be the angle the line makes with the \( x \)-axis.

\[ L = L_{a,\theta} = \{ (x, a + x\tan\theta) : x \in \R \} \]

and teh set of all such lines;

\[ X = \{ L_{a,\theta} : a \in \R \text{ and } \theta \in (\frac{-\pi}{2}, \frac{\pi}{2}) \} \]

With distance function \( \Delta \text { on } X \);

\[ \Delta(L_{a,\theta}, L_{b,\phi}) = \mid a - b \mid + \mid \tan\theta - \tan\phi \mid \]

\qpart

\begin{align*}
\Delta(L_{0,\frac{\pi}{4},L_{1,\frac{-\pi}{4}}}) &= \mid 0 - 1 \mid + \mid \tan\frac{\pi}{4} - \tan\frac{-\pi}{4}\\[8pt]
&= \mid 1 \mid + \mid 1 - (-1) \mid\\[8pt]
&= 1 + 2\\[8pt]
&= 3
\end{align*}

\vspace{2cm}

\qpart

Using \textup{Definition 2.1 Chapter 14 p72};

\qsubpart

\textbf{M1:}

\[ d(a,b) \geq 0 \]
Let \( L_{a,\theta}, L_{b,\phi} \in X \).
\begin{align*}
\Delta(L_{a,\theta}, L_{b,\phi}) &= \mid a - b \mid + \mid \tan\theta - \tan\phi \mid\\[8pt]
&\geq 0 + 0\\[8pt]
&= 0
\end{align*}

Since absolute values are always non-negative, non-negativity holds.

\vspace{1cm}

Let \( L_{a,\theta}, L_{b,\phi} \in X \).
\begin{align*}
\Delta(L_{a,\theta}, L_{b,\phi}) &= \mid a - b \mid + \mid \tan\theta - \tan\phi \mid = 0\\[8pt]
\stext{This holds if and only if both absolute value terms are zero}
&\iff \mid a - b \mid = 0 \text{ and } \mid \tan\theta - \tan\phi \mid = 0\\[8pt]
&\iff a = b \text{ and } \tan\theta = \tan\phi\\[8pt]
&\iff a = b \text{ and } \theta 
= \phi \quad \text{(since } \theta, \phi \in (\frac{-\pi}{2}, \frac{\pi}{2})\text{)}\\[8pt]
&\iff L_{a,\theta} = L_{b,\phi}
\end{align*}

\vspace{1cm}

\qsubpart

\textbf{M2:}

\[ d(a,b) = d(b,a) \]

\begin{align*}
\Delta(L_{a,\theta}, L_{b,\phi}) &= \mid a - b \mid + \mid \tan\theta - \tan\phi \mid\\[8pt]
&= \mid b - a \mid + \mid \tan\phi - \tan\theta \mid\\[8pt]
&= \Delta(L_{b,\phi}, L_{a,\theta})
\end{align*}

Thus, symmetry holds.

\vspace{1cm}

\qsubpart

\textbf{M3:}

\[ d(a,c) \leq d(a,b) + d(b,c) \]

Let \( L_{a,\theta}, L_{b,\phi}, L_{c,\psi} \in X \).
\begin{align*}
\Delta(L_{a,\theta}, L_{c,\psi}) &= \mid a - c \mid + \mid \tan\theta - \tan\psi \mid\\[8pt]
&\leq \mid a - b \mid + \mid b - c \mid + \mid \tan\theta - \tan\phi \mid + \mid \tan\phi - \tan\psi \mid\\[8pt]
&= \Delta(L_{a,\theta}, L_{b,\phi}) + \Delta(L_{b,\phi}, L_{c,\psi})
\end{align*}

Thus, the triangle inequality holds.

\vspace{3cm}

\qpart

Using \textup{Definition 2.6 Chapter 14 p72};

\[ B_{\Delta}(L_{0,0},1) \]

\begin{align*}
B_{\Delta}(L_{0,0},1) &= \{ L_{b,\phi} \in X : \Delta(L_{0,0}, L_{b,\phi}) < 1 \}\\[8pt]
&= \{ L_{b,\phi} \in X : \mid 0 - b \mid + \mid \tan0 - \tan\phi \mid < 1 \}\\[8pt]
&= \{ L_{b,\phi} \in X : \mid b \mid + \mid 0 - \tan\phi \mid < 1 \}\\[8pt]
\stqxt{Applying the triangle inequality in \( \R \) to each absolute value term}
&= \{ L_{b,\phi} \in X : \mid b \mid + \mid \tan\phi \mid < 1 \}
\end{align*}

\vspace{2cm}

\qpart

Is the sequence of lines given by;

\[ (L_{n})_{n=1}^{\infty}, L_{n} = L_{\frac{-1}{n},\frac{\pi}{2}-\frac{1}{n}} \]

\( \Delta-convergent \)?

To be \( \Delta-convergent \), there must exist some line \( L_{b,\phi} \in X \) such that for every \( \epsilon > 0 \), there exists some \( N \in \N \) such that for all \( n \geq N \);
\[ \Delta(L_{n}, L_{b,\phi}) < \epsilon \]
Calculating \( \Delta(L_{n}, L_{b,\phi}) \);
\begin{align*}
\Delta(L_{n}, L_{b,\phi}) &= \Delta(L_{\frac{-1}{n},\frac{\pi}{2}-\frac{1}{n}}, L_{b,\phi})\\[8pt]
&= \mid \frac{-1}{n} - b \mid + \mid \tan(\frac{\pi}{2}-\frac{1}{n}) - \tan\phi \mid\\[8pt]
\end{align*}    

\end{question}

Alternative

\begin{question}

\[ L = \mathbf{L_{a}} ,\theta = \{ (x,a+x\tan\theta):x \in \R \} \]

Where \( a \) is the y- intercept of the line and \( \theta \) is the angle of inclination of the line to the positive x-axis.

\[ \mathbf{X} =  \mathbf{L_{a}} ,\theta : \mathbf{a} \in \R \text{ and } \theta \in (-\pi/2, \pi/2)  \]

and

\[ \Delta(\mathbf{L_{a}},\theta, \mathbf{L_{b}}, \phi) = \mid a-b \mid + \mid \tan{\theta} - \tan{\phi} \mid \]
for \( \mathbf{L_{a}},\theta, \mathbf{L_{b}}, \phi \in \mathbf{X} \)

\qpart

The distance between \( \mathbf{L_{0}},\pi/4 \text{ and } \mathbf{L_{1}},-\pi/4 \);

\begin{align*}
\Delta(\mathbf{L_{0}},\pi/4, \mathbf{L_{1}},-\pi/4) &= \mid 0-1 \mid + \mid \tan{\pi/4} - \tan{-\pi/4} \mid \\[8pt]
&= 1 + \mid 1 - (-1) \mid \\[8pt]
&= 1 + \mid 2 \mid \\[8pt]
&= 1 + 2 = 3
\end{align*}

\qpart

Bt \textup{Definition 2.1,HB Chapter 14 p72},

\qsubpart

To show \( \Delta \) satisfies (M1) non-negativity, we need to show that for any 
\( d^{(2)}(\mathbf{a},\mathbf{b}) \geq 0 \)

\begin{align*}
\Delta(\mathbf{L_{a}},\theta, \mathbf{L_{b}}, \phi) &= \mid a-b \mid + \mid \tan{\theta} - \tan{\phi} \mid \\[8pt]
\stext{As \( \mathbf{a,b} \in \R \) and \( \theta \in (-\pi/2, \pi/2) \), both \( \mid a-b \mid \geq 0 \) and \( \mid \tan{\theta} - \tan{\phi} \mid \geq 0 \)} \\[8pt]
&\geq 0 + 0 \\[8pt]
\end{align*}

Thus, \( \Delta(\mathbf{L_{a}},\theta, \mathbf{L_{b}}, \phi) \geq 0 \) and (M1) is satisfied.

\vspace{2cm}

\qsubpart

To show \( \Delta \) satisfies (M2) symetry, we need to show that for any
\( d^{(2)}(\mathbf{a},\mathbf{b}) = d^{(2)}(\mathbf{b},\mathbf{a}) \)

\begin{align*}
\Delta(\mathbf{L_{a}},\theta, \mathbf{L_{b}}, \phi) &= \mid a-b \mid + \mid \tan{\theta} - \tan{\phi} \mid \\[8pt]
\stext{By the properties of absolute values, HB p8, \( \mid a-b \mid = \mid b-a \mid \) and \( \mid \tan{\theta} - \tan{\phi} \mid = \mid \tan{\phi} - \tan{\theta} \mid \)} \\[8pt]
&= \mid b-a \mid + \mid \tan{\phi} - \tan{\theta} \mid \\[8pt]
&= \Delta(\mathbf{L_{b}},\phi, \mathbf{L_{a}}, \theta)
\end{align*}

Thus, \( \Delta(\mathbf{L_{a}},\theta, \mathbf{L_{b}}, \phi) = \Delta(\mathbf{L_{b}},\phi, \mathbf{L_{a}}, \theta) \) and (M2) is satisfied.

\vspace{2cm}

\qsubpart

To show \( \Delta \) satisfies (M3) triangle inequality, we need to show that for any
\( d^{(2)}(\mathbf{a},\mathbf{c}) \leq d^{(2)}(\mathbf{a},\mathbf{b}) + d^{(2)}(\mathbf{b},\mathbf{c}) \)

\begin{align*}
\Delta(\mathbf{L_{a}},\theta, \mathbf{L_{c}}, \psi) &= \mid a-c \mid + \mid \tan{\theta} - \tan{\psi} \mid \\[8pt]
\stext{By the triangle inequality for real numbers, HB p9, \( \mid a-c \mid \leq \mid a-b \mid + \mid b-c \mid \) and \( \mid \tan{\theta} - \tan{\psi} \mid \leq \mid \tan{\theta} - \tan{\phi} \mid + \mid \tan{\phi} - \tan{\psi} \mid \)} \\[8pt]
&\leq \left( \mid a-b \mid + \mid b-c \mid \right) + \left( \mid \tan{\theta} - \tan{\phi} \mid + \mid \tan{\phi} - \tan{\psi} \mid \right) \\[8pt]
&= \mid a-b \mid + \mid \tan{\theta} - \tan{\phi} \mid + \mid b-c \mid + \mid \tan{\phi} - \tan{\psi} \mid \\[8pt]
&= \Delta(\mathbf{L_{a}},\theta, \mathbf{L_{b}}, \phi) + \Delta(\mathbf{L_{b}},\phi, \mathbf{L_{c}}, \psi)
\end{align*}

Thus, \( \Delta(\mathbf{L_{a}},\theta, \mathbf{L_{c}}, \psi) \leq \Delta(\mathbf{L_{a}},\theta, \mathbf{L_{b}}, \phi) + \Delta(\mathbf{L_{b}},\phi, \mathbf{L_{c}}, \psi) \) and (M3) is satisfied.

\vspace{5cm}

\qpart

\[ \mathbf{B_{\Delta}}(\mathbf{L_{0,0}},1) \]

By \textup{Definition 2.6, HB Chapter 14 p72},

\begin{align*}
\mathbf{B_{\Delta}}(\mathbf{L_{0,0}},1) &= \{ \mathbf{L_{b}}, \theta \in \mathbf{X} : \Delta(\mathbf{L_{0,0}}, \mathbf{L_{b}}, \theta) < 1 \} \\[8pt]
&= \{ \mathbf{L_{b}}, \theta \in \mathbf{X} : \mid 0-b \mid + \mid \tan{0} - \tan{\theta} \mid < 1 \} \\[8pt]
&= \{ \mathbf{L_{b}}, \theta \in \mathbf{X} : \mid b \mid + \mid 0 - \tan{\theta} \mid < 1 \} \\[8pt]
&= \{ \mathbf{L_{b}}, \theta \in \mathbf{X} : \mid b \mid + \mid \tan{\theta} \mid < 1 \}
\end{align*}

\vspace{2cm}

\qpart

The sequence of lines;
\[ (\mathbf{L_{n}})_{n=1}^{\infty} \text{ given by } \mathbf{L_{n}} = \mathbf{L}_{-1/n,\pi/2-1/n} \]
is not \( \Delta-convergent \) because, by \textup{Definition 3.2 HB Chapter 14 p73};

\begin{align*}
\stext{Using \(L_{b,\phi} \in X\) Then}
\Delta(L_n, L_{b,\phi}) &= \left|-\frac{1}{n} - b\right|+ \left|\tan\left(\frac{\pi}{2}-\frac{1}{n}\right) - \tan\phi\right| \\[8pt]
\stext{As \(n \to \infty\), we have \(-1/n \to 0\), so}
\left|-\frac{1}{n} - b\right| \to |b|\\[8pt]
\stext{Also, as \(n \to \infty\), we have}
\tan\left(\frac{\pi}{2}-\frac{1}{n}\right) \to +\infty \\[8pt]
\stext{and since \(\tan\phi\) is finite for all \(\phi \in (-\pi/2,\pi/2)\),}
\left|\tan\left(\frac{\pi}{2}-\frac{1}{n}\right) - \tan\phi\right| \to \infty.
\end{align*}

Hence \(\Delta(L_n, L_{b,\phi}) \to \infty\), and so
\((L_n)\) is not \(\Delta\)-convergent.

\end{question}

\end{document}